\documentclass[11pt]{article}
\usepackage[a4paper,margin=1in]{geometry}
\usepackage{amsmath,amssymb,amsthm,mathtools}
\usepackage{graphicx}
\usepackage{float}
\usepackage[hidelinks]{hyperref}
\usepackage{microtype}

\newtheorem{theorem}{Theorem}
\newtheorem{lemma}[theorem]{Lemma}
\theoremstyle{remark}
\newtheorem{remark}[theorem]{Remark}

\newcommand{\D}{\mathbb D}
\newcommand{\N}{\mathbb N}
\newcommand{\K}{\mathcal K}
\newcommand{\rhoh}{\rho}
\newcommand{\diag}{\operatorname{diag}}

\title{Infinitely many nontrivial root blocks\\in a Parseval operator-orbit frame}
\author{Candidate full solution packet for arXiv:2605.29671}
\date{12 August 2026}

\begin{document}
\maketitle

\begin{center}
\fbox{\parbox{0.91\textwidth}{\textbf{Status.}
Candidate full solution, likely valid, awaiting specialist review. For every
fixed integer \(m\ge2\), we construct a Parseval single-operator orbit frame
whose generator has infinitely many nontrivial root blocks of length \(m\)
and a Riesz basis of generalized eigenvectors.}}
\end{center}

\section{Source question}

Gallardo-Guti\'errez and Partington study frames generated by operator orbits.
In the model-space case, a pair \((A,u)\) generating a frame is similar to
\((S_\theta,k_0^\theta)\), where

\[
 K_\theta=H^2\ominus\theta H^2,
 \qquad S_\theta=P_{K_\theta}S|_{K_\theta}.
\]

After discussing repeated zeros of the inner function \(\theta\), they write:

\begin{quote}
``It does not seem to be known whether generalized eigenvectors can form a
Riesz sequence (if the zeros form a Carleson set and only finitely-many zeros
are repeated, then this will happen), so we do not know whether it is possible
to have infinitely-many nontrivial blocks if there is a frame.''
\end{quote}

This appears on PDF page 17, Section 6.2(2), of arXiv:2605.29671
\cite{source}.

\begin{figure}[H]
 \centering
 \includegraphics[width=0.97\textwidth]{figures/open_question_crop.png}
 \caption{The open question on PDF page 17 of the source paper.}
\end{figure}

\section{Proof intuition}

An interpolating Blaschke product \(b\) has a Riesz basis of normalized
kernels at its simple zeros. Raising \(b\) to a fixed power \(m\) repeats
\emph{every} zero \(m\) times, while the model space retains the transparent
orthogonal decomposition
\[
 K_{b^m}=K_b\oplus bK_b\oplus\cdots\oplus b^{m-1}K_b.
\]
Thus \(\{b^r e_j\}\) is a known Riesz basis. We compare it with the value and
derivative kernels \(\{d_{j,s}:0\le s<m\}\), which are the root vectors of the
backward shift.

The cross-Gram operator is upper triangular in the derivative order: a factor
\(b^r\) has a zero of order \(r\) at every zero of \(b\). Its diagonal blocks
contain the normalized derivatives
\((1-|a_j|^2)b'(a_j)\), uniformly bounded away from zero by Carleson
separation. Hence every diagonal block, and then the entire finite triangular
operator, is invertible. This upgrades Besselness of the derivative kernels
to a Riesz basis.

\section{Preliminaries}

We use the convention that the Hardy-space inner product is linear in its
first variable. For \(a\in\D\), let
\[
 D_{a,s}(z)=\frac{s!z^s}{(1-\overline a z)^{s+1}},
 \qquad \langle f,D_{a,s}\rangle=f^{(s)}(a).
\]

We need one standard sampling fact.

\begin{lemma}[separated derivative sampling]\label{lem:bessel}
Let \(\{a_j\}\subset\D\) be uniformly separated in the pseudohyperbolic
metric. For every fixed \(s\ge0\), the family
\[
 d_{j,s}:=(1-|a_j|^2)^{s+1/2}D_{a_j,s}
\]
is Bessel in \(H^2\).
\end{lemma}

\begin{proof}
For \(s=0\), uniform separation makes
\(\sum_j(1-|a_j|^2)\delta_{a_j}\) a Carleson measure, giving the assertion.
For \(s\ge1\), choose pairwise disjoint pseudohyperbolic disks of one fixed
small radius about the \(a_j\). Subharmonicity on those disks gives
\[
 (1-|a_j|^2)^{2s+1}|f^{(s)}(a_j)|^2
 \le C_s\int_{\Delta_j}|f^{(s)}(z)|^2
                   (1-|z|^2)^{2s-1}\,dA(z).
\]
Summing over the disjoint disks and using the standard weighted Hardy
derivative estimate
\[
 \int_\D |f^{(s)}(z)|^2(1-|z|^2)^{2s-1}\,dA(z)
 \le C'_s\|f\|_{H^2}^2
\]
proves the Bessel inequality. This is the usual derivative form of Hardy
sampling; the \(s=0\) interpolation statement is the
Shapiro--Shields theorem \cite{ShapiroShields}.
\end{proof}

\section{Construction and solution}

\begin{theorem}\label{thm:main}
For every fixed integer \(m\ge2\), there are an infinite-dimensional Hilbert
space \(H\), a contraction \(A\in\mathcal L(H)\), and \(u\in H\) such that:
\begin{enumerate}
 \item \(\{A^n u:n\ge0\}\) is a Parseval frame for \(H\);
 \item \(A\) and \(A^*\) each have a Riesz basis of generalized eigenvectors;
 \item there are infinitely many distinct eigenvalues, and the root block at
 each of them has length exactly \(m\).
\end{enumerate}
In particular, a frame can coexist with infinitely many nontrivial blocks.
\end{theorem}

\begin{proof}
For \(j\ge1\), put
\[
 a_j=\tanh(j+1).
\]
The subtraction identity for the hyperbolic tangent gives
\[
 \rhoh(a_j,a_k)=\left|\frac{a_j-a_k}{1-a_ja_k}\right|
 =\tanh|j-k|.
\]
Moreover \(\sum_j(1-a_j)<\infty\), so these are the zeros of a Blaschke
product \(b\). If
\(p=\prod_{r=1}^{\infty}\tanh r\), then \(p>0\), since
\(\sum_r(1-\tanh r)<\infty\). Therefore
\[
 \inf_j\prod_{k\ne j}\rhoh(a_j,a_k)\ge p^2>0.
\]
Thus \(b\) is an interpolating Blaschke product. The normalized kernels
\[
 e_j(z)=\frac{(1-a_j^2)^{1/2}}{1-a_jz}
\]
form a Riesz basis of \(K_b\) by the Shapiro--Shields theorem
\cite{ShapiroShields}; see also \cite{GMR}.

Set \(\theta=b^m\), \(H=K_\theta\), \(A=S_\theta\), and
\(u=k_0^\theta\). The standard product decomposition of model spaces gives
\[
 K_{b^m}=\bigoplus_{r=0}^{m-1}b^rK_b.
\]
Multiplication by each \(b^r\) is isometric, so
\[
 \mathcal Y=\{y_{j,r}:=b^r e_j:j\ge1,\ 0\le r<m\}
\]
is a Riesz basis of \(K_\theta\).

For \(0\le s<m\), define
\[
 d_{j,s}=(1-a_j^2)^{s+1/2}D_{a_j,s}.
\]
Every \(d_{j,s}\) belongs to \(K_\theta\): if \(h\in H^2\), then
\((b^mh)^{(s)}(a_j)=0\), hence \(d_{j,s}\perp b^mH^2\).
Lemma~\ref{lem:bessel}, applied finitely many times, shows that the combined
family
\(\mathcal D=\{d_{j,s}:j\ge1,0\le s<m\}\) is Bessel.

Consider the cross-Gram operator from \(\mathcal D\) to \(\mathcal Y\), in
blocks indexed by \(0\le r,s<m\). Its entries are
\[
 \langle y_{i,r},d_{j,s}\rangle
 =(1-a_j^2)^{s+1/2}(b^re_i)^{(s)}(a_j).
\]
This vanishes when \(s<r\), so the block matrix is upper triangular. On the
diagonal \(s=r\), Leibniz' rule gives
\begin{align*}
 \langle y_{i,r},d_{j,r}\rangle
 &=r!(1-a_j^2)^{r+1/2}b'(a_j)^r e_i(a_j)\\
 &=r!\beta_j^r\langle e_i,e_j\rangle,
 \qquad \beta_j=(1-a_j^2)b'(a_j).
\end{align*}
For a Blaschke product with simple zeros,
\[
 |\beta_j|=\prod_{k\ne j}\rhoh(a_j,a_k)\ge p^2,
\]
while Schwarz--Pick gives \(|\beta_j|\le1\). Consequently each diagonal
block is a product of the invertible Gram operator of the Riesz basis
\(\{e_j\}\) and an invertible diagonal operator. The off-diagonal blocks are
bounded because both \(\mathcal Y\) and \(\mathcal D\) are Bessel. A finite
upper-triangular block operator with invertible diagonal blocks is invertible.

If \(V_{\mathcal Y}\) and \(V_{\mathcal D}\) are the two synthesis operators,
the cross-Gram operator is \(V_{\mathcal Y}^*V_{\mathcal D}\), up to the
harmless transpose dictated by the inner-product convention. Since
\(V_{\mathcal Y}\) and this cross-Gram operator are invertible,
\(V_{\mathcal D}\) is invertible. Thus \(\mathcal D\) is a Riesz basis of
\(K_\theta\).

Directly from the displayed formula for \(D_{a,s}\),
\[
 (S^*-\overline a I)D_{a,s}=sD_{a,s-1}\quad(s\ge1),
 \qquad (S^*-\overline a I)D_{a,0}=0.
\]
Since \(K_\theta\) is \(S^*\)-invariant, \(S_\theta^*=S^*|_{K_\theta}\).
Hence \(\mathcal D\) is a Riesz basis of root vectors for \(A^*\), with a
length-\(m\) block at every \(\overline{a_j}=a_j\). In the normalized basis,
\[
 (A^*-a_jI)d_{j,s}=s(1-a_j^2)d_{j,s-1},
\]
and all these coefficients are nonzero. Writing \(A^*=VLV^{-1}\) in this
Riesz basis gives \(A=V^{-*}L^*V^*\). The images under \(V^{-*}\) of the
standard coordinate vectors are therefore a Riesz basis of root vectors for
\(A\), with the same block lengths.

It remains to check the frame. For \(f\in K_\theta\), invariance under \(S^*\)
and the reproducing property at zero give
\[
 \langle f,A^nu\rangle
 =\langle (A^*)^nf,k_0^\theta\rangle
 =(S^{*n}f)(0)=\widehat f(n).
\]
Parseval's identity in \(H^2\) now yields
\[
 \sum_{n\ge0}|\langle f,A^nu\rangle|^2
 =\sum_{n\ge0}|\widehat f(n)|^2=\|f\|_{H^2}^2.
\]
Thus the orbit is a Parseval frame, completing the proof.
\end{proof}

\section{Scope, verification, and novelty}

The theorem answers the literal existence question more strongly by allowing
any prescribed uniform finite block length. There is a normalization caveat:
the coefficients \(s(1-a_j^2)\) tend to zero. Thus the result does not claim
a bounded similarity to the direct sum of canonical Jordan matrices having
unit superdiagonal. It does prove infinitely many nontrivial root blocks and
a Riesz basis of generalized eigenvectors, exactly the properties asked for.

The proof is analytic and uses no computer calculation. The principal review
points are the standard derivative-sampling lemma and the orientation of the
finite triangular cross-Gram operator. Both are recorded separately in
\texttt{VERIFICATION.md}.

A bounded primary-source search through 12 August 2026 used the exact arXiv
id and title together with the phrases \emph{generalized eigenvectors},
\emph{Riesz basis}, \emph{model space}, \emph{repeated zeros},
\emph{derivative kernels}, and \emph{infinitely many Jordan blocks}. It found
background results on model-space derivatives but no exact resolution of the
source question. Novelty confidence is moderate; specialist review is
recommended before any priority claim.

\begin{thebibliography}{9}
\bibitem{source}
E. A. Gallardo-Guti\'errez and J. R. Partington,
\emph{Frame constructions associated with operator orbits},
arXiv:2605.29671 (2026).

\bibitem{ShapiroShields}
H. S. Shapiro and A. L. Shields,
\emph{On some interpolation problems for analytic functions},
Amer. J. Math. \textbf{83} (1961), 513--532.

\bibitem{GMR}
S. R. Garcia, J. Mashreghi, and W. T. Ross,
\emph{Introduction to Model Spaces and Their Operators},
Cambridge Studies in Advanced Mathematics 148, Cambridge University Press,
2016.
\end{thebibliography}

\end{document}
